\scanpage{52}
de Poincaré'' rigide-analytique, disant que sur $X^{\mathrm{an}}$, le complexe de De-Rham est bien une résolution du faisceau des fonctions constantes à valeurs dans le corps de base. Notez bien que c'est certainement trivial pour l'analytique p-adique ordinaire, étant une simple question de séries converge\deleted{\unsure{a}}ntes dans \unsure{ce} cas, mais que la question pour la topologie rigide-analyique (qui est \unsure{beamucoup|beaucoup} plus grossière) est plus délicate. Je vous avoue d'ailleurs que j'ai de grands doutes que les choses marchent aussi simplement que ça, i.e. que la cohomologie rigide-analytique à coéfficients constants tout bêtes donne les ``bons'' nombres de Betti. Je crois me rappeler que Tate \unsure{m'avait|n'avait} dit que c'était déjà faux pour le $H^1$ des courbes elliptiques - il faudrait lui demander quel était son argument; apparemment, on trouverait un peu plus que pour la cohomologie étale, mais quand même pas toute la cohomologie. Il reste plausible néanmoins, quoi qu'il en soit, que la topologie rigide-analytique aura son rôle à jouer dans ces questions.

Notez qu'en caractéristique $p>0$, on a un homomorphisme évident de $H^*(X,\underline F_p)$ dans la cohomologie de De-Rham, comme on voit en calculant cette dernière pour la cohomologie étale et non pour la cohomologie de Zariski (ce qui donne le même résultat, puisque les composantes du complexe de De Rham sont quasi-cohérents), et utilisant la suite spectrale en cohomologie étale
\[
H^*(X)\Longleftarrow H^p(X,\underline H^q(\underline\Omega)).
\]
Cet homomorphisme se factorise d'ailleurs à travers $H^*(X,F_p)-H^*(X,\underline O_X)$, ce dernier s'envoyant dans $H^*(X)$ grace à l'homomorphisme de puissance p-ème $f\rightsquigarrow f^p$, induisant un isomorphisme $\unsure{\underline O_X|\underline\Omega_X}\simeq\underline H^0(\underline\Omega)$. Le composé des homomorphismes canoniques $H^{\deleted{\unsure{p}}\inserted{n}}(X,\underline O_X)\to H^n(X)\to H^n(X,\underline O_X)$ (ce dernier
\margin{51}
